%  LaTeX support: latex@mdpi.com 
%  For support, please attach all files needed for compiling as well as the log file, and specify your operating system, LaTeX version, and LaTeX editor.

%=================================================================
\documentclass[Electronics,article,submit,pdftex,moreauthors]{Definitions/mdpi} 
% For posting an early version of this manuscript as a preprint, you may use "preprints" as the journal and change "submit" to "accept". The document class line would be, e.g., \documentclass[preprints,article,accept,moreauthors,pdftex]{mdpi}. This is especially recommended for submission to arXiv, where line numbers should be removed before posting. For preprints.org, the editorial staff will make this change immediately prior to posting.

%--------------------
% Class Options:
%--------------------
%----------
% journal
%----------
% Choose between the following MDPI journals:
% acoustics, actuators, addictions, admsci, adolescents, aerospace, agriculture, agriengineering, agronomy, ai, algorithms, allergies, alloys, analytica, animals, antibiotics, antibodies, antioxidants, applbiosci, appliedchem, appliedmath, applmech, applmicrobiol, applnano, applsci, aquacj, architecture, arts, asc, asi, astronomy, atmosphere, atoms, audiolres, automation, axioms, bacteria, batteries, bdcc, behavsci, beverages, biochem, bioengineering, biologics, biology, biomass, biomechanics, biomed, biomedicines, biomedinformatics, biomimetics, biomolecules, biophysica, biosensors, biotech, birds, bloods, blsf, brainsci, breath, buildings, businesses, cancers, carbon, cardiogenetics, catalysts, cells, ceramics, challenges, chemengineering, chemistry, chemosensors, chemproc, children, chips, cimb, civileng, cleantechnol, climate, clinpract, clockssleep, cmd, coasts, coatings, colloids, colorants, commodities, compounds, computation, computers, condensedmatter, conservation, constrmater, cosmetics, covid, crops, cryptography, crystals, csmf, ctn, curroncol, currophthalmol, cyber, dairy, data, dentistry, dermato, dermatopathology, designs, diabetology, diagnostics, dietetics, digital, disabilities, diseases, diversity, dna, drones, dynamics, earth, ebj, ecologies, econometrics, economies, education, ejihpe, electricity, electrochem, electronicmat, electronics, encyclopedia, endocrines, energies, eng, engproc, ent, entomology, entropy, environments, environsciproc, epidemiologia, epigenomes, est, fermentation, fibers, fintech, fire, fishes, fluids, foods, forecasting, forensicsci, forests, foundations, fractalfract, fuels, futureinternet, futureparasites, futurepharmacol, futurephys, futuretransp, galaxies, games, gases, gastroent, gastrointestdisord, gels, genealogy, genes, geographies, geohazards, geomatics, geosciences, geotechnics, geriatrics, hazardousmatters, healthcare, hearts, hemato, heritage, highthroughput, histories, horticulturae, humanities, humans, hydrobiology, hydrogen, hydrology, hygiene, idr, ijerph, ijfs, ijgi, ijms, ijns, ijtm, ijtpp, immuno, informatics, information, infrastructures, inorganics, insects, instruments, inventions, iot, j, jal, jcdd, jcm, jcp, jcs, jdb, jeta, jfb, jfmk, jimaging, jintelligence, jlpea, jmmp, jmp, jmse, jne, jnt, jof, joitmc, jor, journalmedia, jox, jpm, jrfm, jsan, jtaer, jzbg, kidney, kidneydial, knowledge, land, languages, laws, life, liquids, literature, livers, logics, logistics, lubricants, lymphatics, machines, macromol, magnetism, magnetochemistry, make, marinedrugs, materials, materproc, mathematics, mca, measurements, medicina, medicines, medsci, membranes, merits, metabolites, metals, meteorology, methane, metrology, micro, microarrays, microbiolres, micromachines, microorganisms, microplastics, minerals, mining, modelling, molbank, molecules, mps, msf, mti, muscles, nanoenergyadv, nanomanufacturing, nanomaterials, ncrna, network, neuroglia, neurolint, neurosci, nitrogen, notspecified, nri, nursrep, nutraceuticals, nutrients, obesities, oceans, ohbm, onco, oncopathology, optics, oral, organics, organoids, osteology, oxygen, parasites, parasitologia, particles, pathogens, pathophysiology, pediatrrep, pharmaceuticals, pharmaceutics, pharmacoepidemiology, pharmacy, philosophies, photochem, photonics, phycology, physchem, physics, physiologia, plants, plasma, pollutants, polymers, polysaccharides, poultry, powders, preprints, proceedings, processes, prosthesis, proteomes, psf, psych, psychiatryint, psychoactives, publications, quantumrep, quaternary, qubs, radiation, reactions, recycling, regeneration, religions, remotesensing, reports, reprodmed, resources, rheumato, risks, robotics, ruminants, safety, sci, scipharm, seeds, sensors, separations, sexes, signals, sinusitis, skins, smartcities, sna, societies, socsci, software, soilsystems, solar, solids, sports, standards, stats, stresses, surfaces, surgeries, suschem, sustainability, symmetry, synbio, systems, taxonomy, technologies, telecom, test, textiles, thalassrep, thermo, tomography, tourismhosp, toxics, toxins, transplantology, transportation, traumacare, traumas, tropicalmed, universe, urbansci, uro, vaccines, vehicles, venereology, vetsci, vibration, viruses, vision, waste, water, wem, wevj, wind, women, world, youth, zoonoticdis 

%---------
% article
%---------
% The default type of manuscript is "article", but can be replaced by: 
% abstract, addendum, article, book, bookreview, briefreport, casereport, comment, commentary, communication, conferenceproceedings, correction, conferencereport, entry, expressionofconcern, extendedabstract, datadescriptor, editorial, essay, erratum, hypothesis, interestingimage, obituary, opinion, projectreport, reply, retraction, review, perspective, protocol, shortnote, studyprotocol, systematicreview, supfile, technicalnote, viewpoint, guidelines, registeredreport, tutorial
% supfile = supplementary materials

%----------
% submit
%----------
% The class option "submit" will be changed to "accept" by the Editorial Office when the paper is accepted. This will only make changes to the frontpage (e.g., the logo of the journal will get visible), the headings, and the copyright information. Also, line numbering will be removed. Journal info and pagination for accepted papers will also be assigned by the Editorial Office.

%------------------
% moreauthors
%------------------
% If there is only one author the class option oneauthor should be used. Otherwise use the class option moreauthors.

%---------
% pdftex
%---------
% The option pdftex is for use with pdfLaTeX. If eps figures are used, remove the option pdftex and use LaTeX and dvi2pdf.

%=================================================================
% MDPI internal commands
\firstpage{1} 
\makeatletter 
\setcounter{page}{\@firstpage} 
\makeatother
\pubvolume{1}
\issuenum{1}
\articlenumber{0}
\pubyear{2022}
\copyrightyear{2022}
%\externaleditor{Academic Editor: Firstname Lastname}
\datereceived{} 
%\daterevised{} % Only for the journal Acoustics
\dateaccepted{} 
\datepublished{} 
%\datecorrected{} % Corrected papers include a "Corrected: XXX" date in the original paper.
%\dateretracted{} % Corrected papers include a "Retracted: XXX" date in the original paper.
\hreflink{https://doi.org/} % If needed use \linebreak
%\doinum{}
%------------------------------------------------------------------
% The following line should be uncommented if the LaTeX file is uploaded to arXiv.org
%\pdfoutput=1

%=================================================================
% Add packages and commands here. The following packages are loaded in our class file: fontenc, inputenc, calc, indentfirst, fancyhdr, graphicx, epstopdf, lastpage, ifthen, lineno, float, amsmath, setspace, enumitem, mathpazo, booktabs, titlesec, etoolbox, tabto, xcolor, soul, multirow, microtype, tikz, totcount, changepage, attrib, upgreek, cleveref, amsthm, hyphenat, natbib, hyperref, footmisc, url, geometry, newfloat, caption

%=================================================================
%% Please use the following mathematics environments: Theorem, Lemma, Corollary, Proposition, Characterization, Property, Problem, Example, ExamplesandDefinitions, Hypothesis, Remark, Definition, Notation, Assumption
%% For proofs, please use the proof environment (the amsthm package is loaded by the MDPI class).

%=================================================================
% Full title of the paper (Capitalized)
\Title{Simplex Lattice Design and X-Ray Diffraction for Analysis of Soil Structure: A Case of Cement-Stabilized Compacted Tills Reinforced with Steel Slag and Slaked Lime}

% MDPI internal command: Title for citation in the left column
\TitleCitation{Simplex Lattice Design and X-Ray Diffraction for Analysis of Soil Structure: A Case of Cement-Stabilized Compacted Tills Reinforced with Steel Slag and Slaked Lime}

% Author Orchid ID: enter ID or remove command
\newcommand{\orcidauthorA}{0000-0002-0577-9936} % Add \orcidA{} behind the author's name
\newcommand{\orcidauthorB}{0000-0002-5759-1089} % Add \orcidA{} behind the author's name

% Authors, for the paper (add full first names)
\Author{Per Lindh $^{1,2}$\orcidA{} and Polina Lemenkova $^{3}\orcidB{}$*}

%\longauthorlist{yes}

% MDPI internal command: Authors, for metadata in PDF
\AuthorNames{Per Lindh, Polina Lemenkova}

% MDPI internal command: Authors, for citation in the left column
\AuthorCitation{Lindh, P.; Lemenkova, P.}
% If this is a Chicago style journal: Lastname, Firstname, Firstname Lastname, and Firstname Lastname.

% Affiliations / Addresses (Add [1] after \address if there is only one affiliation.)
\address{%
$^{1}$ \quad Swedish Transport Administration, Malm\"{o}, Sweden; per.a.lindh@gmail.com \\
$^{2}$ \quad Lund University, Lunds Tekniska H\"{o}gskola LTH (Faculty of Engineering), Department of Building and Environmental Technology, Division of Building Materials, Box 118, SE-221-00, Lund, Sweden.\\
$^{3}$ \quad Universit\'{e} Libre de Bruxelles, \'{E}cole polytechnique de Bruxelles (Brussels Faculty of Engineering), Laboratory of Image Synthesis and Analysis, Building L, Campus de Solbosch, ULB -- LISA CP165/57, Avenue Franklin D. Roosevelt 50, B-1050, Brussels, Belgium}

% Contact information of the corresponding author
\corres{Correspondence: polina.lemenkova@ulb.be; Tel.: +32471860459 (P. Lemenkova)}

% Current address and/or shared authorship
% \secondnote{These authors contributed equally to this work.}
% The commands \thirdnote{} till \eighthnote{} are available for further notes

%\simplesumm{} % Simple summary

%\conference{} % An extended version of a conference paper

% Abstract (Do not insert blank lines, i.e. \\) 
\abstract{Evaluating the structure of soil prior to building construction has become more valuable in a large variety of geotechnical and civil engineering applications. To built an effective framework for assessment and modelling stabilized soil specimens, the workflow can include a complex approach of simplex lattice design and X-Ray diffraction for analysis of soil structure. Different from traditional in-situ measurements, we propose a statistical framework for effective decision on binder combination in the process of soil stabilization. Practical aim of this study was to evaluate strength properties of different soil types using ordinary Portland cement, slaked lime and steel slag as pure agents and blended binders. Soil specimens were collected in southern Sweden and include sandy silty tills and clay till (clay content 6-18\%). The pre-processing included mineralogical analysis of mineral composition and soil structure by X-ray diffraction (XRD) and sieve. The specimens were fabricated, compacted, rammed, stabilized by six binder blends and assessed for Uniaxial Compressive Strength (UCS). Moisture Condition Value (MCV) and water content tests were done for compacted soil and shown variation in the MCV values for different binders. The study determined the effects from binder blends on the UCS gain in three types of soil, measured on days 7, 28 and 90th. Positive effects were noted from steel slag/lime blends on UCS gain in sandy silty tills. Steel slag/slaked lime mixed binder performs better compared to pure binders. The effectiveness of simplex lattice design is demonstrated in a series fo ternary diagrams showing soil strength evaluated by adding stabilising agents in different proportions.}

% Keywords
\keyword{steel slag; simplex lattice design; ordinary Portland cement; slaked lime; X-Ray diffraction; compressive strength; soil stabilization; civil engineering} 

% The fields PACS, MSC, and JEL may be left empty or commented out if not applicable
%\PACS{J0101}
%\MSC{}
%\JEL{}

%%%%%%%%%%%%%%%%%%%%%%%%%%%%%%%%%%%%%%%%%%
% Only for the journal Diversity
%\LSID{\url{http://}}

%%%%%%%%%%%%%%%%%%%%%%%%%%%%%%%%%%%%%%%%%%
% Only for the journal Applied Sciences
%\featuredapplication{Authors are encouraged to provide a concise description of the specific application or a potential application of the work. This section is not mandatory.}
%%%%%%%%%%%%%%%%%%%%%%%%%%%%%%%%%%%%%%%%%%

%%%%%%%%%%%%%%%%%%%%%%%%%%%%%%%%%%%%%%%%%%
% Only for the journal Data
%\dataset{DOI number or link to the deposited data set if the data set is published separately. If the data set shall be published as a supplement to this paper, this field will be filled by the journal editors. In this case, please submit the data set as a supplement.}
%\datasetlicense{License under which the data set is made available (CC0, CC-BY, CC-BY-SA, CC-BY-NC, etc.)}

%%%%%%%%%%%%%%%%%%%%%%%%%%%%%%%%%%%%%%%%%%
% Only for the journal Toxins
%\keycontribution{The breakthroughs or highlights of the manuscript. Authors can write one or two sentences to describe the most important part of the paper.}

%%%%%%%%%%%%%%%%%%%%%%%%%%%%%%%%%%%%%%%%%%
% Only for the journal Encyclopedia
%\encyclopediadef{For entry manuscripts only: please provide a brief overview of the entry title instead of an abstract.}

%%%%%%%%%%%%%%%%%%%%%%%%%%%%%%%%%%%%%%%%%%
\begin{document}

%%%%%%%%%%%%%%%%%%%%%%%%%%%%%%%%%%%%%%%%%%
\section{Introduction}

Cementitious materials have positive effects on developing of strength and hardening of clayey soil during stabilization. Cement is the most widely used traditional binder often used for soil stabilization [1–4]. Apart from cement, many other stabilizing agents are being added to binder blends. Replacing the cement by other materials aims to enhance technical performance of stabilization, increase the economic effectivity of the construction process through cheaper materials and reduce environmental impacts from CO2 discharges [5–7]. Besides, the effects from blends fabricated from various binders on soil are more durable and notable [8–10]. The stabilization of soil is owing to the ability of binders to decrease the space between the pores in soil through bonding of particles during the period of curing [9–10]. The decrease of volume in the pores of soil develops over the time of curing and results in the decreased porosity and permeability of soil [8–10]. Thus, cementitious binders increase the compactness and strength of soil, which finally increases its workability and bearing capacity, – the key requirements for soil quality in road construction.	

Many various binders make notable beneficial effects on the development of strength in binder-soil blends in the process of stabilization. For instance, adding slaked lime and steel slag increases Uniaxial Compressive Strength (UCS) of soil, as reported in many studies [11–13]. As for the impact from the mixtures of binders on soil compactness, earlier works shown their positive effects on various types of soil intended for construction of roads [14–16]. Ground Granulated Blast Furnace Slag (GGBS) is a steel waste with pozzolanic properties, received as an industrial by-product from the steelmaking process [17]. 

The fundamental physical properties of metallurgical slag include viscosity and rheological parameters, formed through chemical composition which includes SiO2 (silica), Al2O3 (aluminium oxide) and MgO (magnesium oxide) [18]. This makes steel slag a widely used as binding component in soil stabilization for roadbed structures [19–21], as it has beneficios impacts on pore structure of soil, due to its physico-chemical properties. Because soil as a porous construction material experience shrinkage and decrease in voids with the inclusion of cementitious binders, the compaction and stabilization processes significantly decrease porosity and increase bulk density of soil [22]. This suggests positive effect from additives in soil, which increase its strength and decrease permeability and porosity when added in optimal amounts [23–24]. Another often-studied environmental parameter that is closely related to climate setting is the freeze-thaw durability of soil used as roadbeds, which is actual for countries with cold climate conditions [25–26]. The UCS of soil remarkably increases with the inclusion of binders, and suggests beneficial effects from reinforcing materials on the resistivity of soil to the seasonal climate cycles, such as freeze-thaw processes. 

The chemistry of stabilizing agents should also be considered when dealing with the effects of binders on soil strength [27]. Earlier works reported that the microstructure and geotechnical parameters of soil (e.g., shear modulus, compaction, Atterberg limits, internal friction angle, and maximum dry unit weight) improve along with added hardeners, such as cement pastes or other stabilizing agents, mainly due to the increased bonds between soil particles, which chemically reinforces its inner structure [28–30]. Steel slag, in particular, binds soil particles effectively in contaminated soil polluted by heavy metals, and decreases leaching [31–33]. In its turn, steel mill sludge can be used for reinforcement of clay bricks, to decrease negative environmental consequences from its discharge and to optimise its utilization [34].

The effects from blended binders and their influence on the development of soil strength are not as straightforward. While in general the inclusion of stabilizing agents increases the UCS of soil, the effects may vary individually depending on the types of binders and soil and external factors (curing period and temperature). For example, slaked lime and fly ashes reduce the permeability and improves the strength in the cohesive soil due to pozzolanic reaction from siliceous and aluminous material [35]. The incorporation of cement decreases the porosity in pore microstructure of soil due to the bonding of particles [36–37]. However, the effects from the soil types and different pH values also increase strength capacity due to higher pH values [38–39], as well as with regard to different soil types (fine grained, medium grained or coarse-grained) [40]. Although rates in strength development do generally increase with increasing content of binders in cement-soil mixture, external factors (moisture content and temperature) also play a significant role. Thus, since the effects of pore refinement in soil and compaction act simultaneously during curing period, it is alway necessary to determine individually in each case the effects from binders that reduce permeability of soil, gain in strength and environmental effects (decreased leaching, reduced pH), which might vary along with the amount and the type of binders used for soil stabilization.

The existing literature and case studies on soil stabilization tended to focus on the final stage assessing the gain in strength. It can be the effects from various binders, both traditional [41–43] or novel binders [44–45], and the evaluation of workability of stabilized soil [46]. These methods are effective to evaluate the general stabilization results and to obtain a rough estimate of the soil strength over a curing period. However, a comparative analysis of blended mixtures needs special attention, to provide adequate data regarding the analysis of the effects from particular binders and their ratios, when mixed in proportions, on soil strength. Apart from the traditional UCS testing used to evaluate the strength of stabilized soil, more advanced techniques that allow for accurate and non-destructive measurements of specimens include seismic waves and X-rays [47–48]. One such technique is the propagating the elastic P-waves through the soil mass, where strength is calculated from the increase in wave velocities which directly correlate with the stiffness, hardness and density in microstructure of the porous media [49–50].

Accurate and systematic evaluation of strength in soil stabilized by various binders using blended mixtures in various proportions are lacking, especially those focused on Swedish clayey tills, which are subject to harsh environmental and climate conditions. The objective of this study is a comparative analysis of the effects from various binders and their proportions on soil strength over a period of curing. Besides, the research on influence of binders on Moisture Condition Value (MCV) of soil need to be performed to show their effects with aim to soil improvement. The non-selective use of binders for any type of soil can result in the opposite effects, since as various types of binders are better adjusted to the various soil types, respectively. Therefore, the effects from diverse binders should be tested and evaluated separately using statistical approach for optimization of the workflow. To this end, we performed a series of statistical simplex tests to define the best relationships of soil-binder mixtures with regard to the strength of soil.

The goal of this study was to determine the effects that stabilizing agents (ordinary Portland cement, slaked lime and slag GGBFS) have on the long-term strength gain in 3 types of soil measured on days 7, 28 and 90th. Besides, other geotechnical characteristics were tested for each type of soil stabilized by various binders, respectively: dry density, water content and MCV. In order to perform tests more effectively, we performed statistical optimization of these processes using existing experimental approaches [51–52]. The research has been carried out in the Swedish Geotechnical Institute (SGI) using the available tools and equipment. The methodology followed the existing standards on soil quality testing defined in technical guidances by the Swedish Institute for Standards (SIS). The dynamics in strength and changes of soil structure stabilized by various binders were examined by the UCS tests, moisture and compactness were monitored with Proctor compaction technique during curing period. The resulting changes were visualised, compared and assessed in series of statistical ternary plots, presented below in the article. The correlations between binders and strength in the stabilized soil mixture, the estimated UCS response surfaces for different recipes of soil-binder mixtures and the observed values were detected, plotted and evaluated.
	
%%%%%%%%%%%%%%%%%%%%%%%%%%%%%%%%%%%%%%%%%%
\section{Materials and Methods}

\subsection{Approach}

To evaluate the structural and physicochemical properties of the soil, we applied two types of tests: the Moisture Condition Value (MCV) for compaction characteristics of specimens; and the Uniaxial Compressive Strength (UCS). These two parameters were analysed to evaluated the compaction properties and gained strength of the stabilized soil.

\subsection{Material characteristics}
The soil specimens used in this study were fabricated from the raw material originated from three different trial pits in southern Sweden. The soil included different types A, B and C with varied physio-chemical characteristics, described as follows. The clay content of soil types A, B and C varied from 6 to 18\%. Based on the dominating size of the particles within tested specimens, soils of types A and C are sandy silty tills, while soil type B is clay till. Soil type A had the following characteristics: particle density $\rho _s =2.70 mg/m^3$ ; Optimum Moisture Content (OMC) according to the modified Proctor is 6.1\%; dry density $\rho _d = 2.26 Mg/m^3$; natural water content (WN) =9.8\%; Liquid limit measured by cone penetrometer (WL) is 19.3\%, plasticity index <10. Soil type B had the following characteristics: particle density $\rho _s = 2.71 mg/m^3$; OMC as modified Proctor – 9.2\%; dry density $\rho _d = 2.12 Mg/m^3$; natural water content (WN) =14.6\%; Liquid limit by cone penetrometer (WL) 23.9\%, plasticity index <11.9. Soil type C, included only measurements of the particle density and WN. Thus, particle density $\rho _s$ for soil type C was $2.67 mg/m^3$ and WN =13.0\%; plasticity index <10. The summary of the performed tests on soil types A, B and C is presented in Table 1. 

We estimated the mineral composition and structure of the solid phases of soil types A, B and C using the difference in X-ray diffraction (XRD), to understand its response on addition of binders during stabilization process. The principle consists in the adsorption of binders during reactions with soil, and transport of the materials of stabilizing agents through pores of soil, which performs at interfaces of soil with various physico-chemical properties that are different for types A, B and C in soil components. Soil types A and B contained clay minerals in a mixed layer of smectite (montmorillonite) (Na,Ca)0.33(Al,Mg)2(Si4O10)(OH)2*nH2O, illite (K,H3O)(Al,Mg,Fe)2(Si,Al)4O10[(OH)2,(H2O)], kaolinite (Al2(OH)4Si2O5). Soil type C contains chlorite (ClO-2), illite (K,H3O)(Al,Mg,Fe)2(Si,Al)4O10[(OH)2,(H2O)] and kaolinite (Al2(OH)4Si2O5). 


\subsection{Binders}
Binders can be rated according to their effects on strength development in soil (Table 2). For example, the suggested binders for illite and kaolinite include sand, cement and lime. This explained the choice of lime and cement as binders for stabilization of soil type A and B. The chlorite component in soil type C well responds to treatment by cement, which is why the OPC was selected for soil type C. For these reasons, the following commercially available binders were used for soil treatment: ordinary Portland cement (OPC), ground granulated blast furnace slag (GGBFS) and a slaked lime. 

The blends of standardized binders were used in blending operation, which was performed in a mixing with soil, according to standards of SGI for soil composition and recommended binders, and controlled through records used for ternary plots for UCS estimation. The we selected different types of binders for each type of soil, because these binders react differently with soil particles. The calculated difference in reaction products from 6 different binder recipes, which are as follows (in kg): 1) OPC (100); 2) OPC (50)+GGBFS (50); 3) GGBFS (100); 4) GGBFS (50) + lime (50); 5) lime (100); 6) lime (50) + OPC (50). The soil was treated with these binders or their blends. The total amount of binder used was set to 2.5\% of the soils' dry density, because it enables to obtain an expected maximum UCS of 5 MPa after a curing period of 3 months. A UCS of 5 MPa is a high value for soil stabilization, but was obtained using 2.5\% binder.

All three soil types A, B and C contained material varying from clay to coarse sand particles and therefore could be treated with lime and cement. There are differences in the effects from various binders (Table 3). For instance, OPC quickly starts to react with soil particles just within 2 h after mixing with soil. Most of the reaction is finished in 2 months, Fig. 1. In contrast, the effects from lime are more consequent and similar to a straight line. 

To correspond between the long-term soil performance and workflow effectiveness in the geotechnical works, various manufacturers are producing specific binders with trade names. However, often the exact content and the proportions might not be available, which requires testing binders for soil stabilization. For example, lime and OPC mixed with clay produce reaction products after 18 months. However, lime-clay has more productivity compared to cement-clay with regard to reaction products, which is dependent on the temperature in the stabilized soil. OPC works better for coarse material, while lime is better suitable for fine-grained soil. Again, contrary to lime, OPC does not require external materials to result in a binding blend. The coarse silt type of soil can be stabilized by lime and cement better than by the pure binder.


\subsection{Compaction}
The samples were treated by vibratory surface compaction using the equipment available at SGI to compact the granular material of soil. The 4500 g of each soil specimens were compacted in three layers in a plastic tube with an outer steel tube. The surface of each compacted layer was scratched before the adding of the next layer. All the samples were compacted one hour after mixing with binders which excluded the evaluation of curing period from various binders. After compaction, the specimens were cut to a height of ca. 2.06 m, sealed with paraffin and stored within the plastic tube. 

The specimens did not lost weight during the curing period. Each layer was compacted with a 4-5 kg rammer with 25 blows, which resulted in decrease of 30 cm. To sift out large lumps, the soil passed through a 19 mm sieve and lumps larger than 19 mm were removed. The suitable specimens were then moulded again and placed into a sealed container, one by one. The containers were stored in a climate room at a constant temperature of 20°C and a constant relative humidity (RH) of 85\%. The MCV of the mixture, was determined in one hour of soil compaction, while the UCS tests were taken on days 7, 28 and 90th, respectively. 

\subsection{Moisture Condition Value (MCV)}
The MCV of the soil-binder mixture was performed because the compaction is a key procedure of soil improvement with regard to strength and durability. The advantage of MCV consists in its simplicity and availability to measure the lowest compaction energy needed to achieve the maximum soil density at a specific water content for assessing the fitness of soil for geotechnical road works. The MCV was determined in accordance to the standards EN 13286-46 and the UK Standard 1377 [53–54], adapted for current case. This MCV is based on the repeated compaction of specimen with rammer until the limited compaction is reached. 

The technical procedure is as follows. A free-falling 7 kg and 97 mm-diameter rammer was attached to an automatic release device. A soil sample of 1.5 kg was used with a drop height of 2.50 m. A lightweight disk was placed on top of the soil to prevent extrusion of soil and smearing of the rammer. The upper limit for MCV was set up at 14, to ensure the optimum moisture content for stabilised soil. The dry densities of soil specimens, stabilised by 6 types of binders, described in previous section (2.3 Binders) were compacted by the MCA. 

The results are compared with the dry density of the vibratory-compacted samples for soil types A and B, Fig. 2. The dry density is higher for type A compacted by MCA, compared to the vibratory compacted, Fig. 2. In contrast, soil type B demonstrated the same or lower dry density values for the MCA-compaction, compared to the vibratory compaction. The MCV for soil A is slightly higher than that for type B, due to the techniques of the vibratory compaction which does not produce dry density similar to that in MCA compaction for soil A. 

Nevertheless, the vibratory compaction results in identical dry density as the MCA. The MCV values are below 14 for all tested samples. For the case of MCV <14, the soil is workable and becomes homogeneous and dense with normal compaction equipment.  Each soil type was homogenized, however, small variations in the water content remained between each specimen due to the individual properties of soil, which lied within normal distribution.

Water content in the stabilized soil is an important variable for groundwater discharge and evaluation of soil chemistry. Thus, if the moisture content of a soil is optimum, it can be suitable for construction works. At the same time, due to the variations in the mineralogical content (clay, quartz, etc.), natural water content differs in soil types A, B and C within the same group due to the individual properties, and between the investigated soil types.  

\subsection{Simplex Lattice Design}
The mixture design was adopted from [52] to evaluate the effects from individual binders and their blends using simplex lattice design with 3 binders: OPC, GGBFS and slaked lime. We used six components of recipes for normal simplex lattice design. The scheme presents three single blends of 100\% of a single ingredient, and three mixtures of two of the binders (Fig. 4). Three statistical regression models model the material, depending on the test design. These include linear, quadratic and special cubic models. 
(a) Linear model:
(b) Quadratic model:
(c) Special cubic model:

where $\beta$ is the regression coefficient; $\epsilon$ is the residual error; $x_i$ is the value of factors; $y$ is the dependent variable.

The dependent variable y contains effect $x_1<\ldots>x_3$ (linear model), interactions $x_1x_2$, $x_1x_3$, $x_2x_3$ (quadratic model) and $x_1x_2x_3$ (special cubic model). We included only the significant effects from the quadratic and special cubic regression models. The regression models were tested with the ANOVA to evaluate the significance. For models with not significant ANOVA, each effect was tested for significance and those with a p>0.05 were removed.  The null hypothesis $H_0$ was that the mean of MCV or UCS of the stabilized soil was equal for all binders. This alternative hypothesis $H_1$ is that at least one binder differs from the others. Testing the equality of binders leads to the equation (4): 

\begin{equation}
H_0 = \mu _1 = \mu _2 = <\ldots> = \mu _a
\end{equation}

\begin{equation}
H_1: \mu _i \neq \mu _j
\end{equation}
for at least one pair of $(i, j)$.

\section{Results and Discussion}
Results demonstrate a series of plotted ternary plots showing response surfaces for UCS and MCV (Figs. 5 to 8). The circles on the borders of these figures signify a mixture of single binder, binary blend or a complex mixture. The results of the MCV tests are presented for soil types A and B in Fig. 5. Straight lines representing the contours signify no interaction between the binders as independent variables. In other words, no 2nd-order term is significant (equations 1 and 2 for linear and quadratic models). Response surfaces in Figures 5 to 8 is show the effects from binders on soil stabilization. 

\subsection{MCV}
The response surface equation for binder blends tested for MCV on soil types A and B is presented in equations (6) and (7). Here, C, L and S stand for cement, lime and steel slag, and combinations LS and CLS – for their mixes, correspondently. The numbers represent the amount of binders taken for stabilization.

MCV=13.2C+11.4L+8.9S+3.7CS+2.2LS 		(6) 
MCV=12.4C+10.5L+8.6S+7.4CL+22.5CLS	 (7)

The total amount is adjusted for 100\%, that is, boundary conditions satisfy to the following assumption:  $0 \le C$, $S, L \le 1$ and $C+L+S = 1$. The sum of the independent variables is 1 for all cases. One can derive from equations (6) and (7), that cement has the greatest impact on soil compaction after one hour of treatment time. There are correlations between various binders for all the soil types with those having a positive mark indicating the increase in MCV.

%%%%%%%%%%%%%%%%%%%%%%%%%%%%%%%%%%%%%%%%%%
\section{Discussion}

Authors should discuss the results and how they can be interpreted from the perspective of previous studies and of the working hypotheses. The findings and their implications should be discussed in the broadest context possible. Future research directions may also be highlighted.

%%%%%%%%%%%%%%%%%%%%%%%%%%%%%%%%%%%%%%%%%%
\section{Conclusions}

This section is not mandatory, but can be added to the manuscript if the discussion is unusually long or complex.

%%%%%%%%%%%%%%%%%%%%%%%%%%%%%%%%%%%%%%%%%%
\section{Patents}

This section is not mandatory, but may be added if there are patents resulting from the work reported in this manuscript.

%%%%%%%%%%%%%%%%%%%%%%%%%%%%%%%%%%%%%%%%%%
\vspace{6pt} 

%%%%%%%%%%%%%%%%%%%%%%%%%%%%%%%%%%%%%%%%%%
%% optional
%\supplementary{The following supporting information can be downloaded at:  \linksupplementary{s1}, Figure S1: title; Table S1: title; Video S1: title.}

% Only for the journal Methods and Protocols:
% If you wish to submit a video article, please do so with any other supplementary material.
% \supplementary{The following supporting information can be downloaded at: \linksupplementary{s1}, Figure S1: title; Table S1: title; Video S1: title. A supporting video article is available at doi: link.}

%%%%%%%%%%%%%%%%%%%%%%%%%%%%%%%%%%%%%%%%%%
\authorcontributions{For research articles with several authors, a short paragraph specifying their individual contributions must be provided. The following statements should be used ``Conceptualization, X.X. and Y.Y.; methodology, X.X.; software, X.X.; validation, X.X., Y.Y. and Z.Z.; formal analysis, X.X.; investigation, X.X.; resources, X.X.; data curation, X.X.; writing---original draft preparation, X.X.; writing---review and editing, X.X.; visualization, X.X.; supervision, X.X.; project administration, X.X.; funding acquisition, Y.Y. All authors have read and agreed to the published version of the manuscript.'', please turn to the  \href{http://img.mdpi.org/data/contributor-role-instruction.pdf}{CRediT taxonomy} for the term explanation. Authorship must be limited to those who have contributed substantially to the work~reported.}

\funding{Please add: ``This research received no external funding'' or ``This research was funded by NAME OF FUNDER grant number XXX.'' and  and ``The APC was funded by XXX''. Check carefully that the details given are accurate and use the standard spelling of funding agency names at \url{https://search.crossref.org/funding}, any errors may affect your future funding.}

\institutionalreview{In this section, you should add the Institutional Review Board Statement and approval number, if relevant to your study. You might choose to exclude this statement if the study did not require ethical approval. Please note that the Editorial Office might ask you for further information. Please add “The study was conducted in accordance with the Declaration of Helsinki, and approved by the Institutional Review Board (or Ethics Committee) of NAME OF INSTITUTE (protocol code XXX and date of approval).” for studies involving humans. OR “The animal study protocol was approved by the Institutional Review Board (or Ethics Committee) of NAME OF INSTITUTE (protocol code XXX and date of approval).” for studies involving animals. OR “Ethical review and approval were waived for this study due to REASON (please provide a detailed justification).” OR “Not applicable” for studies not involving humans or animals.}

\informedconsent{Any research article describing a study involving humans should contain this statement. Please add ``Informed consent was obtained from all subjects involved in the study.'' OR ``Patient consent was waived due to REASON (please provide a detailed justification).'' OR ``Not applicable'' for studies not involving humans. You might also choose to exclude this statement if the study did not involve humans.

Written informed consent for publication must be obtained from participating patients who can be identified (including by the patients themselves). Please state ``Written informed consent has been obtained from the patient(s) to publish this paper'' if applicable.}

\dataavailability{In this section, please provide details regarding where data supporting reported results can be found, including links to publicly archived datasets analyzed or generated during the study. Please refer to suggested Data Availability Statements in section ``MDPI Research Data Policies'' at \url{https://www.mdpi.com/ethics}. If the study did not report any data, you might add ``Not applicable'' here.} 

\acknowledgments{In this section you can acknowledge any support given which is not covered by the author contribution or funding sections. This may include administrative and technical support, or donations in kind (e.g., materials used for experiments).}

\conflictsofinterest{Declare conflicts of interest or state ``The authors declare no conflict of interest.'' Authors must identify and declare any personal circumstances or interest that may be perceived as inappropriately influencing the representation or interpretation of reported research results. Any role of the funders in the design of the study; in the collection, analyses or interpretation of data; in the writing of the manuscript; or in the decision to publish the results must be declared in this section. If there is no role, please state ``The funders had no role in the design of the study; in the collection, analyses, or interpretation of data; in the writing of the manuscript; or in the decision to publish the~results''.} 

%%%%%%%%%%%%%%%%%%%%%%%%%%%%%%%%%%%%%%%%%%
%% Optional
\sampleavailability{Samples of the compounds ... are available from the authors.}

%% Only for journal Encyclopedia
%\entrylink{The Link to this entry published on the encyclopedia platform.}

\abbreviations{Abbreviations}{
The following abbreviations are used in this manuscript:\\

\noindent 
\begin{tabular}{@{}ll}
MDPI & Multidisciplinary Digital Publishing Institute\\
DOAJ & Directory of open access journals\\
TLA & Three letter acronym\\
LD & Linear dichroism
\end{tabular}
}

%%%%%%%%%%%%%%%%%%%%%%%%%%%%%%%%%%%%%%%%%%
%% Optional
\appendixtitles{no} % Leave argument "no" if all appendix headings stay EMPTY (then no dot is printed after "Appendix A"). If the appendix sections contain a heading then change the argument to "yes".
\appendixstart
\appendix
\section[\appendixname~\thesection]{}
\subsection[\appendixname~\thesubsection]{}
The appendix is an optional section that can contain details and data supplemental to the main text---for example, explanations of experimental details that would disrupt the flow of the main text but nonetheless remain crucial to understanding and reproducing the research shown; figures of replicates for experiments of which representative data are shown in the main text can be added here if brief, or as Supplementary Data. Mathematical proofs of results not central to the paper can be added as an appendix.

\begin{table}[H] 
\caption{This is a table caption.\label{tab5}}
\newcolumntype{C}{>{\centering\arraybackslash}X}
\begin{tabularx}{\textwidth}{CCC}
\toprule
\textbf{Title 1}	& \textbf{Title 2}	& \textbf{Title 3}\\
\midrule
Entry 1		& Data			& Data\\
Entry 2		& Data			& Data\\
\bottomrule
\end{tabularx}
\end{table}

\section[\appendixname~\thesection]{}
All appendix sections must be cited in the main text. In the appendices, Figures, Tables, etc. should be labeled, starting with ``A''---e.g., Figure A1, Figure A2, etc.

%%%%%%%%%%%%%%%%%%%%%%%%%%%%%%%%%%%%%%%%%%
\begin{adjustwidth}{-\extralength}{0cm}
%\printendnotes[custom] % Un-comment to print a list of endnotes

\reftitle{References}

% Please provide either the correct journal abbreviation (e.g. according to the “List of Title Word Abbreviations” http://www.issn.org/services/online-services/access-to-the-ltwa/) or the full name of the journal.
% Citations and References in Supplementary files are permitted provided that they also appear in the reference list here. 

%=====================================
% References, variant A: external bibliography
%=====================================
%\bibliography{your_external_BibTeX_file}

%=====================================
% References, variant B: internal bibliography
%=====================================
\begin{thebibliography}{999}
% Reference 1
\bibitem[Author1(year)]{ref-journal}
Author~1, T. The title of the cited article. {\em Journal Abbreviation} {\bf 2008}, {\em 10}, 142--149.
% Reference 2
\bibitem[Author2(year)]{ref-book1}
Author~2, L. The title of the cited contribution. In {\em The Book Title}; Editor 1, F., Editor 2, A., Eds.; Publishing House: City, Country, 2007; pp. 32--58.
% Reference 3
\bibitem[Author3(year)]{ref-book2}
Author 1, A.; Author 2, B. \textit{Book Title}, 3rd ed.; Publisher: Publisher Location, Country, 2008; pp. 154--196.
% Reference 4
\bibitem[Author4(year)]{ref-unpublish}
Author 1, A.B.; Author 2, C. Title of Unpublished Work. \textit{Abbreviated Journal Name} year, \textit{phrase indicating stage of publication (submitted; accepted; in press)}.
% Reference 5
\bibitem[Author5(year)]{ref-communication}
Author 1, A.B. (University, City, State, Country); Author 2, C. (Institute, City, State, Country). Personal communication, 2012.
% Reference 6
\bibitem[Author6(year)]{ref-proceeding}
Author 1, A.B.; Author 2, C.D.; Author 3, E.F. Title of presentation. In Proceedings of the Name of the Conference, Location of Conference, Country, Date of Conference (Day Month Year); Abstract Number (optional), Pagination (optional).
% Reference 7
\bibitem[Author7(year)]{ref-thesis}
Author 1, A.B. Title of Thesis. Level of Thesis, Degree-Granting University, Location of University, Date of Completion.
% Reference 8
\bibitem[Author8(year)]{ref-url}
Title of Site. Available online: URL (accessed on Day Month Year).
\end{thebibliography}

% If authors have biography, please use the format below
%\section*{Short Biography of Authors}
%\bio
%{\raisebox{-0.35cm}{\includegraphics[width=3.5cm,height=5.3cm,clip,keepaspectratio]{Definitions/author1.pdf}}}
%{\textbf{Firstname Lastname} Biography of first author}
%
%\bio
%{\raisebox{-0.35cm}{\includegraphics[width=3.5cm,height=5.3cm,clip,keepaspectratio]{Definitions/author2.jpg}}}
%{\textbf{Firstname Lastname} Biography of second author}

% For the MDPI journals use author-date citation, please follow the formatting guidelines on http://www.mdpi.com/authors/references
% To cite two works by the same author: \citeauthor{ref-journal-1a} (\citeyear{ref-journal-1a}, \citeyear{ref-journal-1b}). This produces: Whittaker (1967, 1975)
% To cite two works by the same author with specific pages: \citeauthor{ref-journal-3a} (\citeyear{ref-journal-3a}, p. 328; \citeyear{ref-journal-3b}, p.475). This produces: Wong (1999, p. 328; 2000, p. 475)

%%%%%%%%%%%%%%%%%%%%%%%%%%%%%%%%%%%%%%%%%%
%% for journal Sci
%\reviewreports{\\
%Reviewer 1 comments and authors’ response\\
%Reviewer 2 comments and authors’ response\\
%Reviewer 3 comments and authors’ response
%}
%%%%%%%%%%%%%%%%%%%%%%%%%%%%%%%%%%%%%%%%%%
\end{adjustwidth}
\end{document}

